\chapter[Physical Mechanisms]{Physical Mechanisms for Inhibiting Germination}

Many Fabaceae (legumes), a few Malvaceae, and a few other species use
impervious seed coats as the mechanism for preventing germination before the seed
is dispersed. It is not clear whether the mechanism operates by excluding water or
oxygen or both, but what is clear is thit grinding whole through the seed coat
produces a dramatic effect with the seeds germinating in a few days after puncturing
the seed coats and placing in moist media. With seeds large enough to be held in
pliers, the seed is held against a grinding wheel until the interior of the seed just
shows. Alternatively the seeds can be filed or sawed with a hacksaw until the
endosperm shows. The seeds can also be held in pliers or tweezers and ground
against sandpaper or a file. It was best to grind against the convex side of the seed to
avoid any mechanical damage to the undeveloped radicle. The dramatic effects of
such pretreatment of the seed is shown in Table 9-1 at the end of this Chapter.
The favorable effects on germination induced by various grinding and abrasion
procedures has long been known. This subject is treated in recent books under the
title scarification. This term should be promptly abandoned because producing a scar
on the seed coat surface has no effect. Only the formation of a true water channel and
hole through the seed coat will reliably and reproducibly give immediate germination.
Thompson and Morgan use the words nicking, pricking, and filing in their catalog.

Puncturing has been used herein. Obviously we are searching for a suitable word.
It has often been recommended to soak seeds of this type in hot water or to pour
hot or boiling water over the seeds. An example taken from the Thompson and
Morgan catalog is to immerse seeds of Mimosa pudica in 140 deg. F water for twenty
minutes. I have verified the effectiveness of this procedure. Prof. Roger Koide at
Pennsylvania State University has found that ten minutes in water at 140 is effective
for Abutilon theophrasti. The Thompson and Morgan catalog recommends soaking in
hot water for a number of species, and it can be inferred that these have impervious
seed coats. There are two problems with these hot water treatments. The optimum
conditions vary with each species and have to be determined, and the results are
generally less effective than physically producing the hole in the seed coat as shown
below by the results with Gymnocladus dioica.

The hot water treatments work because the heat causes expansion of the seed
coats. This causes microfissures to form and in effect a hole has been produced in the
seed coat. The extensive studies described below on Gymnocladus dioica are all in
accord with this interpretation. Inferences in the literature that the soaking softens the
seed coat are incorrect, although the seed coats then soften in a day or two due to
enzymes secreted by the activated seed.

Gymnocladus dioica seeds were treated for five and ten minutes at ten degree
temperature intervals ranging from 120 to 180 deg F. Eight seeds were used at each
condition. The number germinated in the five minute treatments were 3, 4, and 4 at
120, 130, and 140. The number germinated in the ten minute treatments were 2, 4,
and 2 at 1 .20, 130, and 140. The seeds were all killed at 150-180. Increasing the time
from five minutes to ten minutes increases the number of seeds killed with no added
stimulation of germination. Samples of eight seeds were then given fifteen second
treatments at 150, 160, and 170. The numbers of seeds germinated were 4, 1, and 1.

Clearly the fifteen seconds at 150 was just enough to expand the seed coat without
cooking the seed inside. However, \textit{none of these treatments gave over 50\%
germination which is less effective than the 100\% germination effected by physically
producing a hole in the seed coat.} All germinations were in two to ten days.
Freezing and thawing have often been recommended for inducing germination
of seeds. In general the effect of freezing and thawing is a myth, and true freezing and
formation of ice crystals within the seed would be fatal. However, temperature
oscillations generally serve to open microfissures. Albizziajulibrissin exemplifies this
effect. Seed subjected to outdoor treatment over one winter gave 24\% germination in
the following April and first week of May showing that microfissures had been opened
in 24\% of the seeds. There was additional germination after a second winter. -There
was no germination in controls held at 40 or 70. In nature it may take several years for
a significant number of seeds to form the requisite microfissures, and it has been
observed in the garden here that Kentucky Coffee Tree (Gymnocladus) seeds do not
germinate for several years after falling. Drying can also induce fissures in the seed
coats, and this was evident in Honey Locust and Kentucky Coffee Tree, Table 9-1.

Many other treatments have been recommended, but all of them suffer in that
they give results that lack uniformity and reproducibility, and they require extensive
study to define optimum conditions. These methods include rubbing between
sandpaper, heating until a few seeds pop, centrifugal devices that throw the seed
against abrasive surfaces, and immersion in concentrated sulfuric acid. This last
procedure is particularly hazardous as was described in Chapter 3.
Although impervious seed coats are characteristic of Fabaceae, many species
in this genus have microfissures already present in varying degrees. Peas and beans
of the vegetable garden germinate readily enough without pretreatment. Hedysarum
cephrolotes, Lathyrus latifolius, Lespedeza bicolor, and Wisteria chinensis germinated
immediately at 70. Baptisia australis, Cercis canadensis, Cercis chinensis, and
Indigofera heteranthera showed some increase in rate of germination with
pretreatment but the minor gain in rate was overbalanced by the increase in rotting
which lowered the overall percent germination (Chapter 20). Coronilla varia, Genista
subcapitata, Lotus corniculatus, and a number of Astragalus and Oxytropis would
probably have benefitted from a hole in the seed coat, but the experiments were not
conducted and modest amounts of germination were achieved without puncturing.

Finally some species with impervious seed coats produce a significant
percentage of seeds that are imperfect in that they already have a miprofissure. These
germinate immediately, but the remaining will not germinate until a hole is produced in
the seed coat. The data on Melilotus alba are an excellent example of this effect.
As might be expected it is species with impervious seed coats that have the
longest life in dry storage. This is discussed in Chapter 13.

The presence of a hard bony seed coat does not indicate the need for
mechanical pretreatment, and in fact a number of very hard and bony seeds have a
water channel through the seed coat. Examples are Cornus Prunus, and Rosa
multiflora. careful examination of these seeds will reveal the cleavage between the
two halves, and producing a hole in the seed coats does not initiate germination.
Conversely a soft pliable seed coat can be an impervious seed coat as in Acer negundo.

The dispersal of seeds with impervious seed coats is interesting. Often the
seed is embedded in a pod containing sweet edible material as in the Kentucky Coffee
Tree pods. Animals carry away these pods and eat the contents thus dispersing the
seed.

There have been repeated reports in the literature that physically hard seed
coats serve as a mechanical constraint on germination. I know of no convincing
evidence that seed coats \textit{significantly} inhibit germination by \textit{mechanically} preventing
expansion of the seed. Three examples from my own work illustrate how seeds
germinate despite very hard seed coats. In stone fruits like many Prunus the seed is
disk shaped with a cleft along the keel. Not only does this cleft allow access by water
and oxygen, but it is a weak point which breaks when the seed starts to expand. In
Gymnocladus dioica the seed coats are very hard and impervious. When a hole is
made in this seed coat, water is rapidly imbibed and the seed secretes powerful
enzymes that soften the seed coat. Thus in 13 days the seed coat has become soft
and rubbery and the seed has expanded to 2-3 times its original volume without
rupturing the seed coat. The radicle has no difficulty in penetrating such a soft seed
coat. Many legumes (Fabaceae) exhibit this phenomenon. In Halesia caroliniana the
seed coat is very hard, but when the seed finally decides to germinate the radicle
bursts through the side in a particularly constructed thin spot and the radicle and
cotyledons develop within a week.

There were several unusual examples of delay mechanisms that could be
classified as physical mechanisms. In Centaurea maculosa the seeds remain in the
recepticle until late into the winter. If removed from the recepticle they germinate
immediately on encountering moisture at 70. Close examination shows that the apex
of the recepticle is constricted, and tightly filled with pappus. This prevents the
entrance of moisture. In the winds of winter these receptacle are blown off, dispersed,
and mechanically broken apart. The seeds falls on the ground md are ready to
germinate as soon as warm weather arives.

Populus, Salix, and Tussilago farfara are a bit of a puzzle. There is no question
that the capsules suddenly open, the seeds rapidly disperse by wind, and the seeds
quickly germinate on encountering moisture. The seeds rapidly die on drying. This
could be regarded as a physical mechanism for delaying germination until dispersal.
Several grasses, Pennisetum villosum for exámplehave the seed in the center
of a.cluster of spiny bristles. These could inhibit contact of the seed with moist ground
and delay germination.

The question of green seed and induced dormancies is taken up here because
the only examples of such effects from my work were with seeds with impervious seed
coats. The horticultural literature contains suggestions that immature (green) seed will
sometimes germinate immediately whereas ripe seed will not. While it is possible that
germination inhibitors could be produced at the very last stages of seed ripening, this
is unlikely and has certainly not been demonstrated. What was found was that seeds
of Sophora japonica hang on the tree in a green state for weeks in late fall. If this seed
is removed and washed, it will germinate immediately. However, if this seed is dry
stored for just three months at 70, the seed coat hardens and becomes impervious.

Now the seed coat must be punctured to get the immediate germination recorded in
Table 9-1. These effects are described in detail under Sophora in Chapter 20. It is
likely that this same effect could be found in other legumes. It is particularly easy to
demonstrate with Sophora japonica because the seed hangs on the tree in a green
but fully sized state for weeks. Attempts to duplicate this effect with Gymnocladus were
unsuccessful as all green seed rotted on contacting moisture.

\begin{table}[h!]
\refstepcounter{table}

\begin{center}
Table \thetable. Dramatic Initiation of Germination from Making a Hole in the Seed Coat
\end{center}

\begin{tabular}{|l|l|l|}
\hline
Species & \multicolumn{2}{|c|}{Germination at 70} \\
        & with pretreatment & without pretreatment \\
\hline
\quad Fabaceae & & \\
Albizzia julibrissin (Mimosa Tree)  & 100\% in 2-4 days (a) & zero in 6 months \\
Colutea arborescens (Bladder Senna) & 100\% in 3-5 days (a) & zero in 6 months \\
Cytisus scoparius (Scotch Broom)    & 100\% in 4-7 days     & zero in 2 months \\
Gleditsia triacanthos (Honey Locust)& 100\% in 3-6 days     & 14\% in 2 months (b) \\
Gymnocladus dioica (Kentucky Coffee T.)& 100\% in 2-4 days  & zero in 6 months (c) \\
Laburnum vossi (Golden Chain Tree)     & 100\% in 5-10 days & zero in 6 months \\
Robinia pseudoacacia (Black Locust T.) & 100\% in 4-6 days  & 10\% in 3 months \\
Sophora japonica (DS seed only)     & 100\% in 3-6 days     & zero in 2 months \\
Tephrosia virginiana                & 100\% in 4-6 days     & zero in one month \\
Thermopsis alpina                   & 100\% in 2-3 days     & zero in one month \\
Vicia americana                     & 100\% in 2-4 days     & 65\% in 3 months(d) \\
\quad Malvaceae & & \\
Abutilon theophrasti                & 100\% in 2-3 days     & 10\% in 2 months(e) \\
Spharalcea parvifolia               & 100\% in 2-3 days     & zero in 4 months \\
\quad Anacardiaceae & & \\
Rhus typhina                        & 100\% in 10-12 days   & zero in 6 months \\
\hline
\end{tabular}


(a) If the pretreated seed is placed in a wet towel at 40 instead of 70, all rotted.

(b) This was increased to 40-50\% if the seed was dry stored at 70 for one year.

(c) Dry storage at 70 for six months gave 25\% germination in 2-4 days.

(d) This was reduced by six months of dry storage. The percent germinations
were 6\% for seed dry stored at 70 and 55\% for seed dry stored at 40. The seeds with
a hole in the seed coat still germinated 100\% in 2-4 days.

(e) Dry storage for six months at 70 gave 25\% germination in 1-8 weeks.
\end{table}

